\section{Conclusion}
\label{sec:conclusion}
Fine-tuning a code model on obfuscated programs roughly doubles its accuracy on them, and we asked
what that gain is made of by evaluating in three regimes rather than one: compositions of the
transforms the model saw, an obfuscator family it never saw, and the same programs queried in the
reverse direction. On seen compositions every adaptation works and the choice between them is worth a
few points. Outside them the ordering inverts. Training one adapter on everything --- the obvious use
of the data --- is the worst arm on the unseen family and loses the untrained direction on most
models; the failures are not malformed answers but the trained answer to the wrong question. Merging
per-transform specialists in weight space, an existing technique applied without modification, is the
the one arm never significantly below the untuned model in any regime on any of eight models --- and
significantly above it on the reverse task on seven --- at a cost of a few
forward points where the training distribution is the right fit.

We take three things from this. Evaluating adaptation in a single view can support the opposite
conclusion to evaluating it in several, and the pairing of directions is a cheap way to add a view to
any behaviour-prediction benchmark. Breadth of training data does not buy breadth of competence; how
the data is combined does. And for a deployment where the obfuscator can change, the recommendation
is the boring one: train specialists, merge them, and deploy one adapter.
